\documentclass[11pt]{article}
\usepackage[margin=1in]{geometry}
\usepackage{amsmath,amssymb,amsthm}
\usepackage{hyperref}
\usepackage{enumitem}

\newtheorem{theorem}{Theorem}
\newtheorem{lemma}[theorem]{Lemma}
\newtheorem{corollary}[theorem]{Corollary}
\theoremstyle{definition}
\newtheorem{definition}[theorem]{Definition}
\theoremstyle{remark}
\newtheorem{remark}[theorem]{Remark}

\usepackage{xurl}
\let\oldus\_
\renewcommand{\_}{\oldus\allowbreak}
\newcommand{\pd}{\operatorname{pd}}

\title{A disproof of Conjecture 00000000504\\[2pt]
\large Projective dimension of edge ideals of $r$-partite $r$-uniform hypergraphs is unbounded}
\date{}

\begin{document}
\sloppy
\maketitle

\section{The conjecture}

Conjecture 00000000504 reads (English version): \emph{``Definition: The length of a minimal free
resolution (projective dimension pd). Conjecture: pd $R/I(G)$ has threshold $2r-1$ for $r$-partite
$r$-uniform hypergraphs (generalizing the Froberg-type theorem pd $\le 3$ for bipartite graphs).''}
The Chinese version says the same: the threshold (the Chinese word means threshold / critical value)
of pd $R/I(G)$ for $r$-partite $r$-uniform hypergraphs is $2r-1$, generalizing the theorem pd
$\le 3$ for bipartite graphs.

\paragraph{Readings.} Here $R=k[x_v : v\in V]$ is the polynomial ring over a field $k$ and $I(G)$
is the edge ideal. The parenthetical (``generalizing ... pd $\le 3$ for bipartite graphs'', and
$2\cdot 2-1=3$) forces the upper-bound sense of ``threshold''. We refute:
\begin{itemize}[nosep]
\item[(A)] $\pd_R R/I(G)\le 2r-1$ for every $r$-partite $r$-uniform hypergraph $G$ (for $r=2$ this
  is the quoted bound $\pd\le 3$ for bipartite graphs);
\item[(B)] the supremum of $\pd_R R/I(G)$ over $r$-partite $r$-uniform hypergraphs equals $2r-1$.
\end{itemize}
Both fail for every field $k$ and every $r\ge 1$. \textbf{Not addressed:} a reading of the form
``$\pd\le 2r-1$ if and only if $G$ has some (unspecified) property'', or a reading in which
``threshold'' refers to an invariant other than pd; the text does not determine such a reading.
A lower-bound reading ($\pd\ge 2r-1$ for all such $G$) is not natural given the parenthetical and is
not formalized.

\section{Definitions (as formalized)}

\begin{definition}
A \emph{hypergraph} on a vertex type $V$ is a finite set $E$ of edges, each a finite subset of
$V$. It is \emph{$r$-uniform} if every edge has exactly $r$ vertices. It is \emph{$r$-partite} if
there is a map $\mathrm{part}:V\to\{0,\dots,r-1\}$ (the parts are the fibres
$V_j=\mathrm{part}^{-1}(j)$) such that every edge meets every part $V_j$ in exactly one vertex.
\end{definition}

\begin{definition}
For a field $k$ and $R=k[x_v:v\in V]$ ($V$ finite), the \emph{edge ideal} is
$I(G)=\bigl(\prod_{v\in e}x_v : e\in E\bigr)\subseteq R$.
\end{definition}

\paragraph{Projective dimension.} $\pd_R M$ is Mathlib's
\texttt{CategoryTheory.projectiveDimension} of $M$ in the category of $R$-modules: the least $n$
such that $\operatorname{Ext}^{i}_R(M,-)=0$ for all $i>n$, with values in
$\mathbb{N}\cup\{\infty\}$ (and $\bot$ for $M=0$). We only prove \emph{lower} bounds on $\pd$.
Any free resolution of $M$ of length $\ell$ shows $\pd_R M\le\ell$, so the length of a minimal free
resolution is at least $\pd_R M$; hence every lower bound below is also a lower bound for the
length of the minimal free resolution named in the conjecture.

\section{The counterexample}

For $m\ge 0$ and $r\ge 1$ let $H_{r,m}$ be the \emph{$r$-uniform matching with $m$ edges}: vertex set
$\{0,\dots,m-1\}\times\{0,\dots,r-1\}$ and edges $e_i=\{(i,0),\dots,(i,r-1)\}$, $i<m$. Distinct
edges are disjoint; $H_{r,m}$ is $r$-uniform, and $r$-partite with parts $V_j=\{(i,j):i<m\}$ (each
$e_i$ meets $V_j$ exactly in $(i,j)$). Its edge ideal is generated by the monomials
$f_i=\prod_{j<r}x_{(i,j)}$, which involve pairwise disjoint sets of variables.

\begin{lemma}[colon step]\label{lem:colon}
Let $J$ be the ideal generated by monomials $x^{d}$, $d\in S$, and let $x^{e}$ be a monomial such
that for every $d\in S$ and every variable $v$, $d_v=0$ or $e_v=0$. If $x^{e}g\in J$ then $g\in J$.
\end{lemma}
\begin{proof}
A polynomial lies in a monomial ideal iff each of its monomials is divisible by a generator. If
$x^{a}$ occurs in $g$, then $x^{a+e}$ occurs in $x^e g$, so $d\le a+e$ for some $d\in S$. At a
variable $v$ with $e_v=0$ this gives $d_v\le a_v$; at a variable with $d_v=0$ it is trivial. So
$d\le a$ and every monomial of $g$ is divisible by a generator.
\end{proof}

\begin{lemma}\label{lem:reg}
If $d_0,\dots,d_{m-1}$ are exponent vectors with pairwise disjoint supports, then
$x^{d_0},\dots,x^{d_{m-1}}$ is a weakly regular sequence in $R$: each $x^{d_i}$ is a
nonzerodivisor on $R/(x^{d_0},\dots,x^{d_{i-1}})$.
\end{lemma}
\begin{proof}
Lemma~\ref{lem:colon} with $S=\{d_0,\dots,d_{i-1}\}$ and $e=d_i$.
\end{proof}

\begin{lemma}\label{lem:pd}
Let $V$ be finite and let $f_1,\dots,f_m\in R$ be a weakly regular sequence of polynomials with zero
constant term. Then $m\le\pd_R R/(f_1,\dots,f_m)$.
\end{lemma}
\begin{proof}
Let $\mathfrak p$ be the kernel of the constant-coefficient map $R\to k$ (a maximal ideal, and
prime because $k$ is a domain). The local ring $R_{\mathfrak p}$ is Noetherian, and since localization
is flat and all $f_i\in\mathfrak p$, the images of $f_1,\dots,f_m$ form a regular sequence of
$R_{\mathfrak p}$. By Mathlib, the quotient of a Noetherian local ring by a regular sequence of
length $m$ has projective dimension exactly $m$. The localization of the $R$-module $R/(f)$ at
$\mathfrak p$ is $R_{\mathfrak p}/(f)R_{\mathfrak p}$, and localization does not increase projective
dimension. Hence $m=\pd_{R_{\mathfrak p}}R_{\mathfrak p}/(f)\le \pd_R R/(f)$.
\end{proof}

\begin{theorem}\label{thm:main}
For every field $k$, every $r\ge 1$ and every $m$, $\pd_R R/I(H_{r,m})\ge m$. Consequently:
\begin{enumerate}[nosep]
\item[(A)] for every $r\ge 1$ the bound $\pd_R R/I(G)\le 2r-1$ fails for the $r$-partite
$r$-uniform hypergraph $G=H_{r,2r}$ (which has $\pd\ge 2r$);
\item[(B)] for every $r\ge 1$ and every $N$ there is an $r$-partite $r$-uniform hypergraph with
$\pd_R R/I(G)>N$ (take $H_{r,N+1}$), so the supremum over the class is infinite, not $2r-1$.
\end{enumerate}
\end{theorem}
\begin{proof}
$I(H_{r,m})=(f_0,\dots,f_{m-1})$ with $f_i=x^{d_i}$, $d_i=\sum_{j<r}\delta_{(i,j)}$; the $d_i$ have
pairwise disjoint supports, and $f_i$ has zero constant term because $r\ge1$. Apply
Lemmas~\ref{lem:reg} and~\ref{lem:pd}.
\end{proof}

For $r=2$: four disjoint edges $x_1y_1,\dots,x_4y_4$ form a bipartite graph with $\pd\ge 4>3$, so
the quoted bound ``pd $\le 3$ for bipartite graphs'' is itself false.

\section{Lean formalization}

File \texttt{lean/Conjecture504/Basic.lean} (Lean 4.33.1, Mathlib v4.33.1, only import
\texttt{Mathlib}), namespace \texttt{Conjecture504}.
\begin{itemize}[nosep]
\item \texttt{Hypergraph}, \texttt{Hypergraph.IsUniform}, \texttt{Hypergraph.IsPartite},
  \texttt{edgeIdeal}: the definitions of Section 2.
\item \texttt{mem\_of\_monomial\_mul\_mem} (Lemma~\ref{lem:colon}),
  \texttt{isWeaklyRegular\_monomials} (Lemma~\ref{lem:reg}),
  \texttt{length\_le\_projectiveDimension} (Lemma~\ref{lem:pd}; uses Mathlib's
  \texttt{IsWeaklyRegular.isRegular\_of\_isLocalization\_of\_mem},
  \texttt{ModuleCat.projectiveDimension\_quotient\_eq\_length},
  \texttt{ModuleCat.projectiveDimension\_le\_projectiveDimension\_of\_isLocalizedModule},
  \texttt{IsLocalizedModule.toLocalizedQuotient'}).
\item \texttt{matching}, \texttt{matching\_isUniform}, \texttt{matching\_isPartite},
  \texttt{edgeIdeal\_matching}, \texttt{matching\_pd\_ge} (the first claim of
  Theorem~\ref{thm:main}).
\item Main theorems:
\begin{verbatim}
theorem conjecture504_upper_bound_false (k : Type u) [Field k] (r : Nat)
    (hr : 0 < r) :
    Not (forall (V : Type) [Fintype V] (G : Hypergraph V),
      G.IsUniform r -> G.IsPartite r ->
      projectiveDimension (ModuleCat.of (MvPolynomial V k)
        (MvPolynomial V k / edgeIdeal k G)) <=
        ((2 * r - 1 : Nat) : WithBot ENat))

theorem conjecture504_unbounded (k : Type u) [Field k] (r : Nat) (hr : 0 < r)
    (N : Nat) :
    exists (V : Type) (_ : Fintype V) (G : Hypergraph V),
      G.IsUniform r /\ G.IsPartite r /\
      (N : WithBot ENat) < projectiveDimension
        (ModuleCat.of (MvPolynomial V k) (MvPolynomial V k / edgeIdeal k G))
\end{verbatim}
(ASCII transcription of the Lean source; the Lean quotient symbol is U+29F8, written 	exttt{/} above.)
\end{itemize}
\paragraph{Axiom audit.} \texttt{lake build} succeeds; \texttt{\#print axioms} for both main
theorems reports exactly \texttt{[propext, Classical.choice, Quot.sound]}. No \texttt{sorry}, no
\texttt{native\_decide}, no added axioms.

\end{document}
